% !TeX root = ../technical_report.tex
\chapter{GPU Virtualisation, NVIDIA Device Plugin, and VRAM Quotas}
\label{ch:gpu_virtualisation}

The deployment of dense Vision Foundation Models (VFMs) at high operational throughput necessitates absolute physical determinism and rigorous computational resource management. Given the constraints of physical bare-metal hardware and the prohibitive costs associated with continuous cloud GPU provisioning, the strategic utilisation of a single unified accelerator for concurrent heterogeneous workloads becomes paramount. This chapter elucidates the architectural implementation of NVIDIA GPU virtualisation, explicit memory budgeting, and latency preservation invariants on a single NVIDIA GeForce RTX 5070 Ti running within a Kubernetes (k3s) environment. The ensuing sections detail the integration of the driver subsystem, the mathematical imperatives of time-slicing versus Multi-Process Service (MPS), the explicit allocation of Video Random Access Memory (VRAM) quotas, and the fail-soft execution fallbacks designed to ensure zero-downtime inference.

\section{NVIDIA Container Toolkit and Driver Subsystem}
\label{sec:container_toolkit_and_driver_subsystem}

The foundational layer of the GPU acceleration stack relies upon the rigorous integration of the \texttt{nvidia-container-toolkit} with the \texttt{containerd} runtime underlying the k3s cluster orchestration on the \texttt{hive-mind} host. The operating system baseline, Ubuntu 26.04.1 LTS (codename \textit{resolute}), provides a highly stable kernel environment. Upon this kernel, the proprietary NVIDIA display driver version 615.71.09 is installed. This driver explicitly reports a CUDA user-mode driver compatibility of version 13.4, thereby establishing the upper bound of the compute API available to user-space applications.

Because the CUDA user-mode driver exhibits strict backward compatibility with older runtime libraries, the containerised workloads are permitted to ship an older CUDA runtime without risking API contract violations. However, the hardware in question, the NVIDIA GeForce RTX 5070 Ti, belongs to the Blackwell architecture generation, denoted by compute capability 12.0. This specific microarchitecture imposes a strict lower bound on the runtime: it dictates that the CUDA runtime must be version 12.8 or newer. Earlier binary compilations of deep learning frameworks, such as PyTorch, completely lack the pre-compiled kernel binaries for compute capability 12.0. Attempting to execute workloads on such deprecated runtimes results in immediate missing-kernel-image aborts, thereby crashing the inference service before initialisation completes. The explicit pinning of the CUDA runtime to version 12.8 within the dependency lock files is therefore an absolute necessity rather than an arbitrary preference.

The orchestration of device assignment within the Kubernetes pods is facilitated by the Container Device Interface (CDI) semantics. The CDI provides a standardised mechanism for container runtimes to inject third-party devices, such as GPUs, into the namespace of a container without requiring elevated privileges or the complete exposure of the host's \texttt{/dev} tree. When a Kubernetes pod requests an \texttt{nvidia.com/gpu} resource, the \texttt{nvidia-device-plugin} generates a highly specific CDI specification JSON document. This document instructs \texttt{containerd} to execute a precise sequence of bind mounts and \texttt{cgroup} device rule modifications, ensuring that the container is granted access solely to the requested GPU control nodes and rendering interfaces, whilst strictly preserving the security boundaries of the host system.

\section{The Single-GPU Multi-Tenancy Challenge and Time-Slicing Architecture}
\label{sec:time_slicing_architecture}

In high-throughput enterprise environments, cloud-native deployments typically rely upon hardware-level partitioning mechanisms to isolate concurrent workloads. Enterprise-grade accelerators, such as the NVIDIA Hopper H100 series, possess hardware Multi-Instance GPU (MIG) capabilities. MIG physically partitions the device memory, the Streaming Multiprocessors (SMs), and the L2 cache, thereby providing strict, hardware-enforced isolation and fault tolerance between tenants.

However, consumer and prosumer graphics cards, including the Ada Lovelace and Blackwell architectures (such as the RTX 5070 Ti), fundamentally lack MIG support. Without further mitigation, requesting a single \texttt{nvidia.com/gpu} resource in a Kubernetes deployment results in an exclusive lock on the entire device. Consequently, attempting to schedule a concurrent CI/CD pipeline, a REST API serving instance, and a batch training job on the same physical host would cause all but one workload to languish perpetually in the \texttt{Pending} state due to resource starvation.

To circumvent this architectural limitation, the TaxonVision-MLOps platform necessitates the explicit implementation of software-based time-slicing via the NVIDIA Kubernetes Device Plugin. The configuration manifest located at \texttt{deploy/k8s/nvidia-time-slicing-config.yaml} partitions the single physical RTX 5070 Ti into four logical, assignable virtual GPU resources.

\begin{listing}[H]
\begin{minted}[frame=lines, framesep=2mm, linenos=true, fontsize=\footnotesize]{yaml}
apiVersion: v1
kind: ConfigMap
metadata:
  name: nvidia-device-plugin-config
  namespace: kube-system
data:
  config.yaml: |
    version: v1
    sharing:
      timeSlicing:
        resources:
          - name: nvidia.com/gpu
            replicas: 4
\end{minted}
\caption{NVIDIA Kubernetes Device Plugin Time-Slicing Configuration specifying four logical replicas.}
\label{lst:nvidia_time_slicing}
\end{listing}

Through this configuration, the device plugin advertises exactly four \texttt{nvidia.com/gpu} resources to the \texttt{kube-scheduler}. This resource scheduling dynamic allows the scheduler to successfully bin-pack up to four discrete pods onto the same physical node, effectively deceiving the Kubernetes control plane into treating the node as if it possessed four independent GPUs.

The underlying mechanism of time-slicing relies heavily upon temporal hardware context switching. The NVIDIA driver multiplexes compute cycles by rapidly swapping the execution context of the Streaming Multiprocessors between the active processes. Whilst time-slicing effectively multiplexes the compute capabilities, it is crucial to understand that it provides absolutely no isolation of memory. All virtual slices fundamentally share the same 16 GB of contiguous GDDR7 VRAM. Therefore, isolation must be stringently enforced at the software layer to prevent Out-Of-Memory (OOM) cascading failures.

\section{Comparative Analysis: Multi-Process Service (MPS) vs Temporal Time-Slicing}
\label{sec:mps_vs_time_slicing}

When addressing the challenge of concurrent execution on consumer GPUs, systems engineers typically deliberate between two distinct methodologies: NVIDIA's Multi-Process Service (MPS) and standard temporal time-slicing.

CUDA MPS operates as a background daemon that intercepts CUDA kernel launches from multiple discrete processes and coalesces them into a single, unified execution context. The primary advantage of MPS lies in its capability to exploit hardware queue multiplexing. By submitting kernels from different processes simultaneously to the GPU, MPS can saturate the Streaming Multiprocessors more efficiently, particularly when the individual workloads are small and under-utilise the available parallel compute resources. MPS effectively allows concurrent processes to share the exact same address space and hardware queues.

However, the TaxonVision-MLOps architecture explicitly rejects MPS in favour of standard temporal time-slicing. The decision is rooted in the paramount requirement for stringent fault isolation guarantees and the mitigation of memory oversubscription risks.

Under MPS, if a rogue training trial -- such as an Optuna hyperparameter optimisation worker exploring a massive batch size -- triggers a memory access violation or an unrecoverable CUDA error, the shared execution context is compromised. This failure cascades, subsequently crashing the background MPS daemon and instantaneously terminating all concurrent workloads, including the mission-critical zero-downtime FastAPI serving container. In a production environment where latency and availability SLAs are strictly enforced, such collateral damage is entirely unacceptable.

Standard temporal time-slicing, conversely, relies upon standard kernel-level temporal context switching. Each process maintains its own distinct CUDA context and independent address space. Although the execution is multiplexed sequentially rather than parallelised concurrently, the fault isolation is robust. If the PyTorch model head training worker encounters an OOM exception or a segmentation fault, the Linux kernel gracefully terminates that specific process without corrupting the CUDA contexts of the co-located ONNX Runtime serving instances. This independent OOM recovery protocol ensures that an unstable CI/CD validation agent cannot inadvertently disrupt the stable production API.

Furthermore, MPS introduces severe memory oversubscription risks. Because the workloads share a context, it becomes exceedingly difficult to strictly partition and enforce memory ceilings per process. Temporal time-slicing, despite sharing the physical VRAM pool, allows for precise programmatic intervention, as each process manages its own allocation limits.

\section{VRAM Quota Allocation and Deterministic Memory Budgeting}
\label{sec:vram_quota_allocation}

The fundamental axiom of time-sliced multi-tenancy on a unified 16 GB memory pool is that the sum of the peak memory allocations of all concurrent processes must never exceed the physical capacity, $M_{\text{total}} \le M_{\text{VRAM}}$. To ensure absolute stability and eliminate the risk of the Linux OOM killer terminating critical services, the TaxonVision-MLOps platform defines a rigorous VRAM budget partitioning matrix across the four logical GPU slots.

The 16.0 GB GDDR7 VRAM is partitioned according to the following explicit policy:

\begin{enumerate}
    \item \textbf{Slot 1: FastAPI / ONNX Runtime Serving Instance.} This slot is allocated a strict VRAM budget of 3.5 GB. The inference pipeline utilises INT8 quantised ONNX checkpoints and caches the feature extraction backbone in evaluation mode. As demonstrated by the empirical memory footprint derived during latency profiling, this 3.5 GB allocation provides an ample buffer for high-throughput batching whilst remaining strictly bounded.
    \item \textbf{Slot 2: Jenkins Ephemeral CI / Feature Validation Agent.} The continuous integration pipeline requires a deterministic environment to validate model convergence and assess feature extraction metrics. This slot is allocated 3.5 GB of VRAM, sufficient to execute unit tests and compile small-scale evaluation graphs.
    \item \textbf{Slot 3: Optuna / PyTorch Model Head Training Worker.} The training pipeline, responsible for evaluating the Class-Balanced Loss on long-tail biological datasets, is inherently memory-intensive due to the requirement to store gradient histories and intermediate activation tensors. This worker is allocated the largest budget, precisely 5.0 GB of VRAM.
    \item \textbf{Slot 4: Elastic Headroom / Active Learning Batch Inference.} The final slot serves as an elastic buffer, allocated 4.0 GB of VRAM. It facilitates asynchronous active learning triage queues and periodic batch re-indexing of the cryptographic exemplar database without starving the real-time inference API.
\end{enumerate}

To mathematically and programmatically enforce these limits within PyTorch, the platform relies upon the \texttt{torch.cuda.set\_per\_process\_memory\_fraction} invariant. Within the execution runner (\texttt{src/taxon\_vision/models/training/runner.py}), the function \texttt{configure\_cuda\_memory\_budget} intercepts the initialisation sequence. For the training worker, a target fraction $f$ is defined. Given $M_{\text{train}} = 5.0\text{ GB}$ and $M_{\text{VRAM}} = 16\text{ GB}$, the fractional allocation evaluates to:
\[
f = \frac{M_{\text{train}}}{M_{\text{VRAM}}} = \frac{5.0}{16.0} = 0.3125
\]
By dynamically configuring \texttt{TAXON\_CUDA\_MEMORY\_FRACTION=0.3125} in the pod specifications, the framework rigorously prevents the PyTorch caching allocator from expanding beyond its designated quota. Any attempt to allocate memory beyond this strict ceiling results in an isolated \texttt{CUDAOutOfMemoryError} within the specific container, preserving the integrity of the adjacent serving slots.

\section{Latency Invariants, Zero-Allocation Tensor Management, and Fail-Soft Execution}
\label{sec:latency_invariants}

The multiplexing of computational cycles inherently introduces temporal latency jitter due to the microsecond overheads of hardware context switching. Maintaining sub-25 millisecond p95 inference latency SLAs under heavy multi-tenant load requires rigorous structural optimisations in the hot execution path.

The primary optimisation strategy is the enforcement of zero-allocation tensor memory management within the \texttt{src/taxon\_vision/service/inference.py} module. Dynamic memory allocation during the forward pass triggers the garbage collector, causing catastrophic latency spikes. To mitigate this, the architecture utilises the \texttt{@lru\_cache(maxsize=1)} decorator to strictly cache the instantiation of the \texttt{TaxonClassifier}, the \texttt{ConformalPredictionEngine}, and the \texttt{EnergyOODDetector}. Furthermore, the deterministic image preprocessing transform pipeline is pre-compiled into a static \texttt{torchvision.transforms.Compose} instance. By preventing the re-instantiation of these complex objects and ensuring that input tensors are fed into pre-allocated memory bindings, the inference module guarantees a zero-allocation hot path. This invariant successfully dampens the context-switching latency variance, maintaining highly deterministic execution times despite the presence of adjacent training workloads.

Despite the robustness of the time-slicing and memory budgeting protocols, architectural resilience demands a comprehensive fail-soft mechanism to guarantee continuous operational availability. The \texttt{src/taxon\_vision/inference/onnx\_engine.py} module implements a highly deliberate fallback topology. During the initialisation of the \texttt{ONNXInferenceEngine}, if the CUDA execution provider fails to acquire the necessary hardware context -- whether due to temporary VRAM exhaustion, device plugin desynchronisation, or driver faults -- the session initialisation gracefully degrades to the \texttt{CPUExecutionProvider}.

\begin{listing}[H]
\begin{minted}[frame=lines, framesep=2mm, linenos=true, fontsize=\footnotesize]{python}
import onnxruntime as ort

class ONNXInferenceEngine:
    def __init__(self, model_path: str) -> None:
        opts = ort.SessionOptions()
        opts.graph_optimization_level = ort.GraphOptimizationLevel.ORT_ENABLE_ALL
        # Deliberate fallback to CPUExecutionProvider ensuring zero-downtime availability
        self.session = ort.InferenceSession(
            model_path,
            opts,
            providers=["CPUExecutionProvider"]
        )
\end{minted}
\caption{ONNX Runtime engine instantiation demonstrating the fail-soft fallback configuration.}
\label{lst:onnx_cpu_fallback}
\end{listing}

While execution on the CPU via the \texttt{CPUExecutionProvider} intrinsically increases inference latency and degrades the p95 throughput profile, this intentional fallback ensures that the FastAPI operator interface remains responsive. In a distributed MLOps environment handling rare biological observations, gracefully delayed execution is overwhelmingly preferable to complete systemic failure and HTTP 503 Service Unavailable errors. This resilient design pattern epitomises the core philosophy of the TaxonVision architecture: deterministic resource management underpinned by robust, automated fault tolerance.
